\section{Conclusion}

\label{sec:conclusion}

In conclusion, this study aims to contribute to the growing literature on media-government dynamics by providing causal evidence on how governments respond to critical media coverage, particularly in the context of advertising spending. By focusing on the case of Mexico, where the interplay between media and government has historically been shaped by clientelistic relationships and economic dependence, this research will seek to uncover whether critical reporting leads to reductions in advertising spending or increases for favorable coverage. Understanding this relationship is critical not only for assessing the health of media independence in Mexico but also for drawing broader implications about the potential for governments to manipulate public discourse through economic incentives and censorship